\section*{الفصل \ref{ch:b3:seismology-imaging} --- علم الزلازل والتصوير الأرضي}

\begin{solution}{exo:b3:seismology-imaging:1}
(أ) $\sin i_0 = pv_0/R = 500\times5.80/6371 = 0.4552$، ومنه $i_0 = \ang{27.1}$.
(ب) عند \qty{771}{km}، $r = \qty{5600}{km}$ و$r/v = 5600/11.066 =
\qty{506.1}{s.rad^{-1}}$؛ عند \qty{871}{km}، $r = \qty{5500}{km}$ و$r/v =
5500/11.245 = \qty{489.1}{s.rad^{-1}}$. ينعطف الشعاع حيث $r/v = p = 500$:
$z_t = 771 + 100\times(506.1 - 500)/(506.1 - 489.1) = 771 + 36 =
\qty{807}{km}$.
\end{solution}

\begin{solution}{exo:b3:seismology-imaging:2}
(أ) $(455.5 - 369.6)/10 = \qty{8.59}{s}$ لكل درجة،
و$8.59\times180/\pi = \qty{492.2}{s.rad^{-1}}$. (ب) درجة القوس الواحدة تساوي
$6371\pi/180 = \qty{111.19}{km}$، ومنه تمسح الجبهة الموجية السطح بسرعة
$111.19/8.59 = \qty{12.94}{km.s^{-1}}$، أسرع بكثير من أيّ صخر عند
السطح: يصل الشعاع بميل حادّ. (ج) عند نقطة الانعطاف $v = r_t/p =
5450/492.2 = \qty{11.07}{km.s^{-1}}$، وهي سرعة موجة P في الوشاح السفلي.
\end{solution}

\begin{solution}{exo:b3:seismology-imaging:3}
(أ) $\mu_1 = \rho_1\beta_1^2 = 2700\times3500^2 = \qty{3.31e10}{Pa}$
(\qty{33.1}{GPa}); $\mu_2 = 3300\times4500^2 = \qty{6.68e10}{Pa}$
(\qty{66.8}{GPa}). (ب) $\lambda = cT = 3.79\times20 = \qty{75.8}{km}$، أي 2.17
ضعف سمك الطبقة: تصل الموجة بعمق كبير داخل نصف الفضاء، وهذا ما
يفسّر لماذا تتجاوز سرعتها $\beta_1$. (ج) $4000/3.384 = \qty{1182}{s}$
و$4000/4.333 = \qty{923}{s}$: تصل الموجات ذات الفترة 20 ثانية بتأخّر \qty{259}{s} أي
\qty{4.3}{min}، عن الموجات ذات الفترة 100 ثانية.
\end{solution}

\begin{solution}{exo:b3:seismology-imaging:4}
(أ) ${}_0S_2$: $1/(53.89\times60) = \qty{3.093e-4}{Hz} = \qty{0.3093}{mHz}$؛
${}_0T_2$: $1/(44.01\times60) = \qty{0.3787}{mHz}$. (ب) $1440/53.89 = 26.7$
و$1440/44.01 = 32.7$ تذبذبًا في اليوم. (ج) الفرق هو
$\qty{0.0694}{mHz}$، ومقلوبه \qty{14400}{s}: يلزم زمن تسجيل لا يقل عن \qty{4.0}{h}،
وفي الممارسة يكون أطول بكثير، لأنّ القمم تتّسع بفعل التوهين
وتنشطر.
\end{solution}

\begin{solution}{exo:b3:seismology-imaging:5}
(أ) الوتر بين نقطتين من الكرة تفصل بينهما الزاوية $\Delta$
عند المركز يبلغ طوله $2R\sin(\Delta/2)$، ويُقطع بالسرعة $v$.
(ب) عند \ang{30}: $2\times6371\times\sin\ang{15}/10.0 = \qty{329.8}{s}$، مقابل
\qty{369.6}{s}؛ عند \ang{90}: $2\times6371\times\sin\ang{45}/10.0 =
\qty{901.0}{s}$، مقابل \qty{779.7}{s}. (ج) ينعطف الشعاع المتّجه إلى \ang{30} في
الوشاح العلوي، حيث تبلغ السرعة الحقيقية نحو 8 إلى \qty{10}{km.s^{-1}}، أدنى من
القيمة المنتظمة \qty{10}{km.s^{-1}}؛ أمّا الشعاع المتّجه إلى \ang{90} فيقضي معظم مساره في
الوشاح السفلي بسرعة 11 إلى \qty{13.7}{km.s^{-1}}، وينحني نحوه، فتكون
الأرض الحقيقية أسرع.
\end{solution}

\begin{solution}{exo:b3:seismology-imaging:6}
(أ) بالتربيع، $(2 - \xi)^4 = 16(1 - \xi/3)(1 - \xi)$. بالنشر:
$16 - 32\xi + 24\xi^2 - 8\xi^3 + \xi^4 = 16 - \tfrac{64}{3}\xi + \tfrac{16}{3}\xi^2$,
ومنه $\xi^4 - 8\xi^3 + \tfrac{56}{3}\xi^2 - \tfrac{32}{3}\xi = 0$، أي
$\xi\,(3\xi^3 - 24\xi^2 + 56\xi - 32) = 0$،
و$3\xi^3 - 24\xi^2 + 56\xi - 32 = (\xi - 4)(3\xi^2 - 12\xi + 8)$. الجذر
$\xi = 0$ لا يمثّل موجة. (ب) جذور المعادلة التربيعية هي $2 \pm 2/\sqrt3$،
أي 3.155 و0.8453؛ لا يقع سوى $\xi = 0.8453$ في $(0, 1)$ (و$\xi = 4$
لا يقع فيها)، ويحقّق المعادلة غير المربّعة، إذ يساوي كلّ من الطرفين $4/3$.
إذن $c/\beta = \sqrt{0.8453} = 0.9194$. (ج) $c = 0.9194\times3.46 =
\qty{3.181}{km.s^{-1}}$؛ على مسافة \qty{1000}{km}، $314.4 - 289.0 = \qty{25.3}{s}$
بعد موجة S.
\end{solution}

\begin{solution}{exo:b3:seismology-imaging:7}
صفوف $\mathsf G$ هي $(1, x_i)$، و$\mathsf C_d^{-1} = \operatorname{diag}
(100, 100, 25, 25)$ (\unit{s^{-2}}). تكون المعادلات الناظمية
\[
  \begin{pmatrix} 250 & 23750 \\ 23750 & \num{2812500} \end{pmatrix}
  \begin{pmatrix} a \\ b \end{pmatrix}
  = \begin{pmatrix} 3665 \\ 418125 \end{pmatrix},
\]
ومحدِّدها \num{1.390625e8}. ومنه $b = (250\times418125 -
23750\times3665)/\num{1.390625e8} = \qty{0.12575}{s.km^{-1}}$
و$a = (2812500\times3665 - 23750\times418125)/\num{1.390625e8} =
\qty{2.713}{s}$. التباين المشترك هو $\mathsf N^{-1}$: $\sigma_b =
\sqrt{250/\num{1.390625e8}} = \qty{0.00134}{s.km^{-1}}$, $\sigma_a =
\sqrt{2812500/\num{1.390625e8}} = \qty{0.142}{s}$. السرعة هي
$1/b = \qty{7.95}{km.s^{-1}}$، والخطأ $\sigma_b/b^2 =
\qty{0.08}{km.s^{-1}}$: وهي سرعة في أعلى الوشاح. البواقي هي $-0.001$،
$0.011$، $-0.076$ و$\qty{0.036}{s}$، و$\chi^2 = 0.19$ عند $4 - 2 = 2$
درجتَي حرية: أصغر ممّا هو متوقَّع، ومنه تكون الأخطاء المذكورة، إن صحّ القول،
مبالغًا فيها؛ والملاءمة مقبولة.
\end{solution}

\begin{solution}{exo:b3:seismology-imaging:8}
(أ) $\mathsf G = (30\;\;40)$ صفٌّ واحد: $\mathsf G^{\mathsf T}\mathsf G$ مصفوفةٌ
$2\times2$ من الرتبة واحد، ومنه فهي مفردة؛ ويمكن إضافة أيّ نموذج يحقّق $30m_1 + 40m_2 = 0$
دون تغيير المعطى. (ب) $\mathsf G\mathsf G^{\mathsf T} =
900 + 1600 = 2500$، ومنه $\mathbf m = (30, 40)\times0.50/2500 =
(0.006, 0.008)\,\unit{s.km^{-1}}$. (ج) $\mathsf R = \frac{1}{2500}
\begin{pmatrix} 900 & 1200 \\ 1200 & 1600 \end{pmatrix} =
\begin{pmatrix} 0.36 & 0.48 \\ 0.48 & 0.64 \end{pmatrix}$، أثرها 1: معطى
واحد يحلّ توليفة واحدة من الوسائط. (د) $\mathsf R\,(0.01, 0) =
(0.0036, 0.0048)\,\unit{s.km^{-1}}$: يتوزّع الشذوذ على طول الشعاع، أكثر
في الخلية التي يكون فيها الشعاع أطول.
\end{solution}

\begin{solution}{exo:b3:seismology-imaging:9}
(أ) بالتقريب الخطّي: $\delta t = -(\delta v/v)L/v = -0.03\times200/9.0 =
\qty{-0.667}{s}$؛ وبالضبط، $L/(1.03v) - L/v = 200/9.27 - 200/9.0 =
\qty{-0.647}{s}$. (ب) يبلغ المسار داخل اللوح المندس $100/\cos\ang{40} =
\qty{130.5}{km}$: $-0.03\times130.5/9.0 = \qty{-0.435}{s}$ بالتقريب الخطّي
و$\qty{-0.422}{s}$ بالضبط. (ج) القيمة الدقيقة هي القيمة الخطّية مقسومةً على
$1.03$: خطأ قدره \qty{3}{\percent}، من مرتبة $\delta v/v$ نفسه، كما
هو متوقَّع من نظرية من الرتبة الأولى.
\end{solution}

\begin{solution}{exo:b3:seismology-imaging:10}
(أ) يوجد النمط التوافقي ما دام جزؤه في نصف الفضاء يتلاشى، $c < \beta_2$؛ وعند
قطعه $c = \beta_2$، ومنه $q_2 = 0$ و$\tan(kHq_1) = 0$، مع $kHq_1$ في
$[n\pi, n\pi + \pi/2)$: $kHq_1 = n\pi$. (ب) مع $k = \omega/\beta_2$
و$q_1 = \sqrt{\beta_2^2/\beta_1^2 - 1} = \sqrt{\beta_2^2 - \beta_1^2}/\beta_1$,
$\omega_nH\sqrt{\beta_2^2 - \beta_1^2}/(\beta_1\beta_2) = n\pi$، وهي
الصيغة. (ج) $\sqrt{4.5^2 - 3.5^2} = \sqrt8 = 2.828$؛ $\omega_1 =
\pi\times3.5\times4.5/(35\times2.828) = \qty{0.4998}{rad.s^{-1}}$، أي فترة
$2\pi/0.4998 = \qty{12.57}{s}$: عند الفترات الأطول لا يوجد إلّا النمط
الأساسي. (د) من أجل $n = 0$ يكون تردّد القطع صفرًا: يوجد النمط
الأساسي عند جميع الفترات. عند الفترات القصيرة يتجه $kH \to \infty$ بينما $kHq_1$ يبقى
أقلّ من $\pi/2$، ومنه $q_1 \to 0$ و$c \to \beta_1$.
\end{solution}

\begin{solution}{exo:b3:seismology-imaging:11}
(أ) $\mathsf G^{\mathsf T}(\mathsf G\mathsf G^{\mathsf T} + \varepsilon^2\mathsf I)
= (\mathsf G^{\mathsf T}\mathsf G + \varepsilon^2\mathsf I)\mathsf G^{\mathsf T}$;
نضرب من اليسار في $(\mathsf G^{\mathsf T}\mathsf G + \varepsilon^2
\mathsf I)^{-1}$ ومن اليمين في $(\mathsf G\mathsf G^{\mathsf T} +
\varepsilon^2\mathsf I)^{-1}$، كلاهما قابل للعكس من أجل $\varepsilon > 0$.
(ب) هنا $\mathsf G\mathsf G^{\mathsf T} = \|\mathbf g\|^2 = 2500$، وهو عدد،
ومنه $\mathsf R_\varepsilon = \mathsf G^{\mathsf T}\mathsf G/(2500 + \varepsilon^2)$،
أثرها $2500/(2500 + \varepsilon^2)$. (ج) المعطى المتنبَّأ به هو
$\mathbf g\cdot\mathbf m_\varepsilon = 0.50\times2500/(2500 + \varepsilon^2)$
والباقي $0.50\,\varepsilon^2/(2500 + \varepsilon^2)$. من أجل
$\varepsilon^2 = 0$: الأثر 1، والباقي 0؛ ومن أجل 2500: الأثر 0.5، والباقي
\qty{0.25}{s}؛ ومن أجل \num{10000}: الأثر 0.2، والباقي \qty{0.40}{s}. يعطي المزيد من التخميد
نموذجًا أصغر، واستبانةً أضعف، وعدم تطابق أكبر.
\end{solution}

\begin{solution}{exo:b3:seismology-imaging:12}
(أ) من التمرين 5، $T = (2R/v)\sin(\Delta/2)$، ومنه $\mathrm dT/\mathrm d\Delta
= (R/v)\cos(\Delta/2)$. الوتر الذي يصنع الزاوية $\Delta$ عند المركز يمثّل
قاعدةَ مثلث متساوي الساقين زاويةُ رأسه $\Delta$، فهو يلتقي نصف القطر
عند المصدر بالزاوية $i_0 = (\pi - \Delta)/2$، و$p = R\sin i_0/v =
(R/v)\cos(\Delta/2)$: وهي النتيجة نفسها. (ب) نقطة الانعطاف هي منتصف
الوتر، على بُعد $R\cos(\Delta_1/2)$ من المركز؛ أو بعبارة أخرى
$r_1 = p_1v$. (ج) هنا $p(\Delta)/p_1 = \cos(\Delta/2)/\cos(\Delta_1/2)$،
وتعطي الصيغة $r_1 = R\exp\bigl(-\ln(1/\cos(\Delta_1/2))\bigr) =
R\cos(\Delta_1/2)$، كما في (ب)؛ والسرعة هي $r_1/p_1 = R\cos(\Delta_1/2)\big/
\bigl((R/v)\cos(\Delta_1/2)\bigr) = v$.
\end{solution}

\begin{solution}{pb:b3:seismology-imaging:1}
\textbf{1.} $i_0 = \arcsin 0.8 = \ang{53.13}$ و$p = 0.8/5.0 =
\qty{0.16}{s.km^{-1}}$.

\textbf{2.} ينعطف الشعاع حيث $v = 1/p = \qty{6.25}{km.s^{-1}}$، أي عند
$z = (6.25 - 5.0)/0.05 = \qty{25}{km}$.

\textbf{3.} بتفاضل $\sin i = pv$ على طول الشعاع، $\cos i\,\mathrm di/
\mathrm ds = pg\,\mathrm dz/\mathrm ds = pg\cos i$، ومنه $\mathrm di/\mathrm ds =
pg$: يدور الاتجاه بمعدّل ثابت، وهو انحناء دائرة نصف قطرها
$1/(pg) = 1/(0.16\times0.05) = \qty{125}{km}$. عند نقطة الانعطاف يكون الشعاع
أفقيًّا على عمق \qty{25}{km}، ويقع المركز رأسيًّا على بُعد \qty{125}{km}
فوقه: عند $z = \qty{-100}{km}$، أي على ارتفاع $v_0/g$ فوق السطح، حيث
ستنعدم السرعة المستقرأة.

\textbf{4.} عند المصدر يكون نصف قطر الدائرة عموديًّا على الشعاع،
فيصنع زاوية $i_0$ مع الأفق، ويقع المركز
على بُعد $125\cos i_0$ أفقيًّا؛ وبالتناظر
$X = 2\times125\times\cos i_0 = 2\times125\times0.6 = \qty{150}{km}$، أو
$X = 2\sqrt{1 - p^2v_0^2}/(pg)$.

\textbf{5.} $\mathrm dT = \mathrm ds/v$ مع $\mathrm ds = \mathrm di/(pg)$
و$v = \sin i/p$: $\mathrm dT = \mathrm di/(g\sin i)$. ثم
\[
  T = \frac2g\int_{i_0}^{\pi/2}\frac{\mathrm di}{\sin i} = \frac2g\ln\cot\frac{i_0}2 =
  \frac2g\operatorname{arccosh}\frac1{pv_0} .
\]
هنا $\cot(\ang{26.57}) = 2$، ومنه
$T = 40\ln2 = \qty{27.73}{s}$.

\textbf{6.} من أجل $p = 0.15$: $pv_0 = 0.75$، $\sqrt{1 - 0.5625} = 0.6614$،
$X = 2\times0.6614/0.0075 = \qty{176.4}{km}$ و$T = 40\operatorname{arccosh}
(1.3333) = 40\times0.7954 = \qty{31.81}{s}$. ثم $\Delta T/\Delta X =
4.08/26.4 = \qty{0.155}{s.km^{-1}}$، بين 0.15 و0.16، كما
تقتضي $\mathrm dT/\mathrm dX = p$.

\textbf{7.} $\mu_1 = 2700\times3500^2 = \qty{33.1}{GPa}$
و$\mu_2 = 3300\times4500^2 = \qty{66.8}{GPa}$.

\textbf{8.} $k = 2\pi/(cT) = 2\pi/(3.99\times25) = \qty{0.06299}{rad.km^{-1}}$;
$q_1 = \sqrt{(3.99/3.5)^2 - 1} = 0.5474$; $q_2 = \sqrt{1 - (3.99/4.5)^2} = 0.4624$.

\textbf{9.} $\tan(kHq_1) = 66.825\times0.4624/(33.075\times0.5474) = 1.707$؛ يأخذ
النمط الأساسي $kHq_1 = \arctan1.707 = 1.041$، ومنه
$H = 1.041/(0.06299\times0.5474) = \qty{30.2}{km}$.

\textbf{10.} عند \qty{40}{s}: $k = 2\pi/170 = 0.03696$، $q_1 = 0.6888$،
$q_2 = 0.3287$، $\tan(kHq_1) = 0.964$، $kHq_1 = 0.7671$
و$H = 0.7671/(0.03696\times0.6888) = \qty{30.1}{km}$. تتّفق القيمتان ضمن
دقة السرعات الطورية: سُمك القشرة نحو \qty{30}{km}.

\textbf{11.} $\omega = 0.2513$ و$\qty{0.1571}{rad.s^{-1}}$، $k = 0.06299$
و$\qty{0.03696}{rad.km^{-1}}$: $U \approx 0.0942/0.02603 = \qty{3.62}{km.s^{-1}}$,
أصغر من السرعتين الطوريتين كلتيهما، كما يجب أن يكون حين تزداد $c$ مع
الدور.

\textbf{12.} $\lambda = 3.99\times25 = \qty{99.8}{km}$ و$4.25\times40 =
\qty{170}{km}$. يسبر النمط الأساسي عمقًا يعادل جزءًا
من طول موجته: تقع موجة \qty{25}{s} في القشرة في معظمها، وتصل موجة \qty{40}{s}
بعمق كبير في الوشاح، فتكون سرعتها أقرب إلى $\beta_2$. ويُعدّ تحديدُ
$H$ نفسه بدورين مختلفين تحقّقًا من ملاءمة نموذج الطبقتين.

\textbf{13.} مع أعمدة الخلايا 1 إلى 4 وصفوف A إلى F، بوحدات
\qty{1000}{km} ($\sqrt2\times100/1000 = 0.1414$):
\[
  \mathsf G = \begin{pmatrix}
  0.1 & 0.1 & 0 & 0 \\ 0 & 0 & 0.1 & 0.1 \\ 0.1 & 0 & 0.1 & 0 \\
  0 & 0.1 & 0 & 0.1 \\ 0.1414 & 0 & 0 & 0.1414 \\ 0 & 0.1414 & 0.1414 & 0
  \end{pmatrix}.
\]

\textbf{14.} $\mathsf C_d^{-1} = \operatorname{diag}(400, 400, 400, 400, 100,
100)$. يسهم شعاع الصف أو العمود بمقدار $0.1^2\times400 = 4$، والشعاع القطري
بمقدار $0.1414^2\times100 = 2$. تعبر الخلية 1 الشعاعان A و C بمقدار 4 لكلٍّ منهما، والشعاع E بمقدار 2: 10. والخليتان
1 و2 تتشاركان الشعاع \mbox{\textdir TLT A: 4}؛ والخليتان 1 و3 تتشاركان \mbox{\textdir TLT C: 4}؛ والخليتان 1 و4 تتشاركان \mbox{\textdir TLT E: 2}؛
وهكذا بالتناظر، وهو ما يعطي الصفوف الأربعة المذكورة.

\textbf{15.} يسهم شعاع الصف أو العمود بمقدار $0.1\times400\,\delta t = 40\,
\delta t$ في كلّ من خلاياه، والشعاع القطري $0.1414\times100\,\delta t =
14.14\,\delta t$. الخلية 1: $40(0.16 - 0.36) + 14.14\times(-0.15) = -10.12$؛ الخلية
2: $40(0.16 + 0.11) + 14.14\times(-0.27) = 6.98$؛ الخلية 3: $40(-0.48 - 0.36) +
14.14\times(-0.27) = -37.42$؛ الخلية 4: $40(-0.48 + 0.11) + 14.14\times(-0.15) =
-16.92$.

\textbf{16.} لأجل $(1,1,1,1)$، مجموع كلّ صفّ من $\mathsf N$ يساوي 20. ولأجل
$(1,-1,-1,1)$، يعطي الصفّ 1 القيمة $10 - 4 - 4 + 2 = 4$، وتتّبع الصفوف الأخرى
النمط نفسه. ولأجل $(1,1,-1,-1)$، فإنّ $10 + 4 - 4 - 2 = 8$، وأمّا لأجل $(1,-1,1,-1)$،
فإنّ $10 - 4 + 4 - 2 = 8$. المتجهات متعامدة ونظيمها 2، ومنه
$\mathsf N^{-1} = \sum_k \mathbf u_k\mathbf u_k^{\mathsf T}/(4\lambda_k)$؛
عنصرها $(1,1)$ يساوي $\tfrac14(\tfrac1{20} + \tfrac14 + \tfrac18 + \tfrac18) =
\tfrac{11}{80}$، وعنصره $(1,2)$ يساوي $\tfrac14(\tfrac1{20} - \tfrac14 + \tfrac18
- \tfrac18) = -\tfrac4{80}$، وعنصره $(1,4)$ يساوي $\tfrac14(\tfrac1{20} +
\tfrac14 - \tfrac18 - \tfrac18) = \tfrac1{80}$، وهكذا.

\textbf{17.} $\hat{\mathbf m} = \mathsf N^{-1}\mathbf b$:
$\hat m_1 = (11\times(-10.121) - 4\times6.982 - 4\times(-37.418) - 16.921)/80
= -0.081$؛ وبالمثل $\hat m_2 = 1.844$، $\hat m_3 = -3.706$
و$\hat m_4 = \qty{-0.931}{ms.km^{-1}}$.

 \textbf{18.} $\delta v/v_0 = -v_0\,\delta s = -8.0\times10^{-3}\,\hat m$، وبالنسبة
المئوية $-0.8\,\hat m$: الخلية 1 $+0.07$، الخلية 2 $-1.48$، الخلية 3 $+2.96$، الخلية 4
$\qty{+0.75}{\percent}$.

 \textbf{19.} كلّ عنصر قطري في $\mathsf N^{-1}$ يساوي $11/80$: لكلّ خلية
الخطأ المعياري $\sqrt{0.1375} = \qty{0.371}{ms.km^{-1}}$، أي
$\qty{0.30}{\percent}$. أمّا ترابط الخليتين 1 و2 فيساوي $-4/11 = -0.36$:
إذ تتشاركان الشعاع A، فيمكن أن يُقايَض التباطؤ المُدخَل في إحداهما بالتباطؤ في الأخرى.
أمّا الخليتان 1 و4، اللتان لا تتشاركان إلّا الشعاع القطري E الأقلّ دقّة، فترتبطان ارتباطًا ضعيفًا
وموجبًا، $+1/11 = +0.09$.

 \textbf{20.} البواقي المتوقَّعة هي $0.176$، $-0.464$، $-0.379$، $0.091$،
$-0.143$ و$\qty{-0.263}{s}$؛ والبواقي بعد الملاءمة هي $-0.016$،
$-0.016$، $0.019$، $0.019$، $-0.007$ و$\qty{-0.007}{s}$،
و$\chi^2 = 0.50$ بدرجات حرية $6 - 4 = 2$: متّسقة مع الأخطاء المذكورة،
وإن كان ثمّة انحراف فهو نحو الأصغر.

 \textbf{21.} المتجهات الذاتية هي متجهات $\mathsf N$ والقيم الذاتية
$\lambda/(\lambda + 4)$: $20/24 = 5/6$، $4/8 = 1/2$، $8/12 = 2/3$ مرّتين. كما في
السؤال 16، $R_{11} = \tfrac14(\tfrac56 + \tfrac12 + \tfrac23 + \tfrac23) =
\tfrac23$; $R_{12} = \tfrac14(\tfrac56 - \tfrac12 + \tfrac23 - \tfrac23) =
\tfrac1{12}$; $R_{13} = \tfrac14(\tfrac56 - \tfrac12 - \tfrac23 + \tfrac23) =
\tfrac1{12}$; $R_{14} = \tfrac14(\tfrac56 + \tfrac12 - \tfrac23 - \tfrac23) = 0$.

 \textbf{22.} نُسقِط $\mathbf b$ على المتجهات الذاتية: $\mathbf b\cdot\mathbf u_k/4
= -14.37$، $0.849$، $12.80$ و$-9.40$ للقيم الذاتية 20، 4، 8 و8.
يقسم التقدير المخمَّد هذه القيم على $\lambda + 4$ بدلًا من $\lambda$؛ وبالنسبة إلى
الخلية 3 (المركِّبات $1, -1, -1, 1$)،
$m_3 = -14.37/24 - 0.849/8 - 12.80/12 - 9.40/12 = \qty{-2.555}{ms.km^{-1}}$,
أي شذوذ في السرعة مقداره $\qty{+2.04}{\percent}$. يسحب التخميد كلّ مركِّبة
نحو الصفر، وأشدّها تأثّرًا المركِّبات ضعيفة القيد.

 \textbf{23.} النموذج غير المخمَّد، الذي يلائم المعطيات ضمن أخطائها
($\chi^2 = 0.50$ بدرجتي حرية)، في حين أنّ النموذج المخمَّد يخالفها
بشدّة ($15.8$). لا يلزم التخميد إلّا حين تكون $\mathsf N$ منفردة أو قريبة
من ذلك؛ وهنا قيمتها الذاتية الأصغر، 4، ليست صغيرة، فالتقدير مستقرّ
وأخطاؤه معتدلة.

 \textbf{24.} الصفّ 3 هو $(1/12, 0, 2/3, 1/12)$: التقدير المخمَّد للخلية 3 يساوي
ثلثي قيمتها الحقيقية مضافًا إليه جزء من اثني عشر من الخليتين 1 و4، وهما جارتاها
على الشعاع C والشعاع B، ولا شيء من الخلية 2.

 \textbf{25.} الخلية 3، وهي الخلية السفلى على اليسار، سريعة بنحو
\qty{3}{\percent}؛ والخلية 2 بطيئة. اللوح المندس البارد الكثيف من الغلاف الصخري المحيطي
الذي يندسّ تحت القارة أسرع من الوشاح المحيط به، لذا تكون الخلية
السريعة، العميقة والواقعة على جانب المحيط من K، هي اللوح المندس؛ أمّا الخلية البطيئة فوقه
على الجانب الآخر فتقع حيث يكون إسفين الوشاح ساخنًا.

 \textbf{26.} من السؤالين 18 و19، شذوذ السرعة الموجية P لخلية اللوح المندس هو
\textbf{\qty[text-series-to-math=true, uncertainty-mode=separate]{3.0 \pm 0.3}{\percent}}.
\end{solution}
